\begin{abstract}
Modern artificial intelligence (AI) training and inference workloads exhibit severe heterogeneity in computational demands, memory footprints, and communication patterns. Existing cluster orchestrators (e.g., Kubernetes, Slurm) often rely on coarse-grained static GPU allocations, leading to hardware underutilization, high tail latencies under bursty workload arrivals, and coarse fault recovery overheads. We present \textbf{DRIGS} (Distributed Resource and Intelligent GPU Scheduling), a lightweight, modular, GPU-aware compute runtime designed for scheduling, executing, monitoring, and recovering heterogeneous AI workloads (inference, distributed training, multi-GPU) across multi-GPU hardware. DRIGS decouples core scheduling logic behind abstract protocols, introducing non-blocking control plane orchestration, dynamic worker node auto-registration, SQLite ACID state persistence, and an automated NVML orphan process sweeper.

Benchmarked on dual NVIDIA Tesla T4 GPU hardware, DRIGS's \texttt{PriorityScheduler} reduces P95 queue wait time by \textbf{60.09\%} ($2.740$\,s vs.\ FIFO $6.866$\,s) under $N=15$ jobs/trial with uniform inter-arrival intervals ($\mathcal{U}(0.5,2.0)$\,s), averaged over 5 trials. For multi-GPU workloads, \texttt{TopologyAwareScheduler} achieves a \textbf{22.19\%} higher topology-quality score than random placement on a simulated 8-GPU cluster. In fault tolerance evaluations, DRIGS achieves control-plane rescheduling in \textbf{0.496\,ms}, with complete physical end-to-end wall-clock recovery following a hard \texttt{SIGKILL} in \textbf{129.13\,ms}. At scale ($N=1{,}000$ jobs), DRIGS maintains \textbf{81.7\,jobs/sec} SQLite control-plane enqueue throughput with scheduler decision latencies under \textbf{17\,ms}. At 95\% VRAM saturation, \texttt{MemoryAwareScheduler} enforces admission throttling (job completion rate: 40.0\%) to prevent GPU memory over-commitment; no CUDA OOM exception occurs on admitted jobs. DRIGS is an open-source, modular GPU scheduling research runtime designed for lightweight deployment and experimental evaluation.\footnote{Source code: \url{https://github.com/Omdeepb69/DRIGS}}
\end{abstract}

\noindent\textbf{Keywords:} GPU Scheduling, Distributed AI Workloads, Resource Management, Fault Recovery, Topology Awareness, High-Performance Computing.
